%% Generated by tools/make_handout.py -- edit config/tasks_localization.json and regenerate.
\documentclass[a4paper,10pt]{article}
\usepackage[utf8]{inputenc}
\usepackage[T1]{fontenc}
\usepackage{lmodern}
\usepackage[textwidth=166mm,top=20mm,bottom=18mm]{geometry}
\usepackage{array}
\usepackage{booktabs}
\usepackage{longtable}
\usepackage{textcomp}
\setlength{\parindent}{0pt}
\setlength{\parskip}{0.6ex}
\newcommand{\pts}[1]{\textsc{#1 pts}}
\pagestyle{plain}
\begin{document}


\begin{center}\large\bfseries Intelligente Robotik --- Praktikum\end{center}
\vspace{-6pt}\begin{center}Localising a mecanum robot: Monte-Carlo localisation and ICP scan matching\end{center}
\vspace{2pt}
\noindent\begin{minipage}[t]{0.60\textwidth}\raggedright
\textbf{Sheets L1---L4}          (4 tasks)

Total 130 pts, pass from 50\%. One sheet per task; hand in one protocol per person (max.\ 2 pages) --- the code is shared, the explanation is not. Every number below has to be reproducible from a command you wrote down.
\end{minipage}\hfill\begin{minipage}[t]{0.38\textwidth}
\begin{tabular}{|l|p{3.6cm}|}
\hline
Name & \rule{0pt}{1.4em} \\
\hline
Date & \rule{0pt}{1.4em} \\
\hline
Group & \rule{0pt}{1.4em} \\
\hline
Seed & \rule{0pt}{1.4em} \\
\hline
\end{tabular}
\end{minipage}
\vspace{4pt}

\vspace{2pt}\noindent\textbf{How the numbers are produced}
The graded run is one number and takes about 36 s. One recording of the same drive turns every question after it into a three-second replay, so the sweep tables below are a loop over \texttt{mcl\_report.py}, not four lab sessions.
\vspace{2pt}{\small
\begin{verbatim}
./tools/run_lab.sh grade --task mcl_production --controller student/mcl_template.py --headless
OHM_RECORD=/tmp/l1.jsonl ./tools/run_lab.sh grade --task mcl_production \
        --controller tools/record_scans.py --headless            # one drive, everything it saw
./tools/mcl_report.py /tmp/l1.jsonl --particles 250 --stride 1 --sigma-z 0.5 --timing
\end{verbatim}
}
\vspace{6pt}

\section{L1 --- Where am I? Monte-Carlo localization against the map}

\noindent\hfill\pts{30}\hspace{1em} hall: \texttt{production}\hspace{1em}\texttt{task: mcl\_production}

\vspace{2pt}\noindent\textbf{Task}

The grader drives; you estimate. This time the improvement is measured against your \textbf{odometry}, not against GPS --- the wheels are the sensor that is already fused into the pose you start with, and beating them is the whole claim of the exercise. One scan every twentieth of a second, 360 beams, and a map you did not build: draw N samples from the belief you were given, move them by the odometry, weight each by how well its predicted beams fit the walls, and resample when the weights stop being spread out. Report the mean and the 1\(\sigma\) on \texttt{/<robot>/kf/pose} --- the \(\sigma\) is graded (NEES), so a cloud that is wide enough to hide in is not a passing answer.

\vspace{2pt}\noindent\textbf{What is graded}

\begin{tabular}{|>{\raggedright\hyphenpenalty 500\arraybackslash}p{9.20cm}|>{\raggedright\hyphenpenalty 500\arraybackslash}p{6.00cm}|}
\hline
\textbf{quantity} & \textbf{target (from the task file)} \\
\hline
accuracy (RMSE over the graded window) & \(\le\) 50 mm \\
improvement over raw odometry & \(\ge\) 5\(\times\) \\
worst single error & \(\le\) 250 mm \\
rate of kf/pose & \(\ge\) 5 Hz \\
wall contacts & \(\le\) 0 \\
NEES (\(\sigma\) honest) & [0.05--5] \\
\hline
\end{tabular}

\vspace{2pt}\noindent\textbf{Protocol} --- one table per item. Four blank rows each; add rows if you need them, and write the command that produced a number under the table.

\vspace{4pt}\noindent\textbf{C1.1} \emph{Particles are moved by the odometry's own delta (rot1, translate, rot2), not teleported to the odometry pose}

\begin{tabular}{|>{\raggedright\hyphenpenalty 500\arraybackslash}p{2.87cm}|>{\raggedright\hyphenpenalty 500\arraybackslash}p{2.28cm}|>{\raggedright\hyphenpenalty 500\arraybackslash}p{2.87cm}|>{\raggedright\hyphenpenalty 500\arraybackslash}p{2.28cm}|>{\raggedright\hyphenpenalty 500\arraybackslash}p{3.70cm}|}
\hline
\textbf{question / item} & \textbf{method} & \textbf{measured} & \textbf{target} & \textbf{explanation of the deviation} \\
\hline
\rule{0pt}{2.4em} & \rule{0pt}{2.4em} & \rule{0pt}{2.4em} & \rule{0pt}{2.4em} & \rule{0pt}{2.4em} \\
\rule{0pt}{2.4em} & \rule{0pt}{2.4em} & \rule{0pt}{2.4em} & \rule{0pt}{2.4em} & \rule{0pt}{2.4em} \\
\rule{0pt}{2.4em} & \rule{0pt}{2.4em} & \rule{0pt}{2.4em} & \rule{0pt}{2.4em} & \rule{0pt}{2.4em} \\
\rule{0pt}{2.4em} & \rule{0pt}{2.4em} & \rule{0pt}{2.4em} & \rule{0pt}{2.4em} & \rule{0pt}{2.4em} \\
\hline
\end{tabular}

\vspace{4pt}\noindent\textbf{C2.1} \emph{Weight of a particle from the map: predicted beam end points compared against the distance field}

\begin{tabular}{|>{\raggedright\hyphenpenalty 500\arraybackslash}p{2.87cm}|>{\raggedright\hyphenpenalty 500\arraybackslash}p{2.28cm}|>{\raggedright\hyphenpenalty 500\arraybackslash}p{2.87cm}|>{\raggedright\hyphenpenalty 500\arraybackslash}p{2.28cm}|>{\raggedright\hyphenpenalty 500\arraybackslash}p{3.70cm}|}
\hline
\textbf{question / item} & \textbf{method} & \textbf{measured} & \textbf{target} & \textbf{explanation of the deviation} \\
\hline
\rule{0pt}{2.4em} & \rule{0pt}{2.4em} & \rule{0pt}{2.4em} & \rule{0pt}{2.4em} & \rule{0pt}{2.4em} \\
\rule{0pt}{2.4em} & \rule{0pt}{2.4em} & \rule{0pt}{2.4em} & \rule{0pt}{2.4em} & \rule{0pt}{2.4em} \\
\rule{0pt}{2.4em} & \rule{0pt}{2.4em} & \rule{0pt}{2.4em} & \rule{0pt}{2.4em} & \rule{0pt}{2.4em} \\
\rule{0pt}{2.4em} & \rule{0pt}{2.4em} & \rule{0pt}{2.4em} & \rule{0pt}{2.4em} & \rule{0pt}{2.4em} \\
\hline
\end{tabular}

\vspace{4pt}\noindent\textbf{C3.1} \emph{N\_eff = 1/\(\Sigma\)w\textsuperscript{2} as the trigger for resampling, and the resampling does not happen on every scan}

\begin{tabular}{|>{\raggedright\hyphenpenalty 500\arraybackslash}p{2.87cm}|>{\raggedright\hyphenpenalty 500\arraybackslash}p{2.28cm}|>{\raggedright\hyphenpenalty 500\arraybackslash}p{2.87cm}|>{\raggedright\hyphenpenalty 500\arraybackslash}p{2.28cm}|>{\raggedright\hyphenpenalty 500\arraybackslash}p{3.70cm}|}
\hline
\textbf{question / item} & \textbf{method} & \textbf{measured} & \textbf{target} & \textbf{explanation of the deviation} \\
\hline
\rule{0pt}{2.4em} & \rule{0pt}{2.4em} & \rule{0pt}{2.4em} & \rule{0pt}{2.4em} & \rule{0pt}{2.4em} \\
\rule{0pt}{2.4em} & \rule{0pt}{2.4em} & \rule{0pt}{2.4em} & \rule{0pt}{2.4em} & \rule{0pt}{2.4em} \\
\rule{0pt}{2.4em} & \rule{0pt}{2.4em} & \rule{0pt}{2.4em} & \rule{0pt}{2.4em} & \rule{0pt}{2.4em} \\
\rule{0pt}{2.4em} & \rule{0pt}{2.4em} & \rule{0pt}{2.4em} & \rule{0pt}{2.4em} & \rule{0pt}{2.4em} \\
\hline
\end{tabular}

\vspace{4pt}\noindent\textbf{C4.1} \emph{kf/pose with sx/sy from the spread of the cloud, not a constant}

\begin{tabular}{|>{\raggedright\hyphenpenalty 500\arraybackslash}p{2.87cm}|>{\raggedright\hyphenpenalty 500\arraybackslash}p{2.28cm}|>{\raggedright\hyphenpenalty 500\arraybackslash}p{2.87cm}|>{\raggedright\hyphenpenalty 500\arraybackslash}p{2.28cm}|>{\raggedright\hyphenpenalty 500\arraybackslash}p{3.70cm}|}
\hline
\textbf{question / item} & \textbf{method} & \textbf{measured} & \textbf{target} & \textbf{explanation of the deviation} \\
\hline
\rule{0pt}{2.4em} & \rule{0pt}{2.4em} & \rule{0pt}{2.4em} & \rule{0pt}{2.4em} & \rule{0pt}{2.4em} \\
\rule{0pt}{2.4em} & \rule{0pt}{2.4em} & \rule{0pt}{2.4em} & \rule{0pt}{2.4em} & \rule{0pt}{2.4em} \\
\rule{0pt}{2.4em} & \rule{0pt}{2.4em} & \rule{0pt}{2.4em} & \rule{0pt}{2.4em} & \rule{0pt}{2.4em} \\
\rule{0pt}{2.4em} & \rule{0pt}{2.4em} & \rule{0pt}{2.4em} & \rule{0pt}{2.4em} & \rule{0pt}{2.4em} \\
\hline
\end{tabular}

\vspace{4pt}\noindent\textbf{Pointers}

\begin{itemize}\itemsep1pt

\item Beam angles: \texttt{angle\_min\ +\ i\(\cdot\)angle\_increment\ +\ theta} of the particle --- a beam that keeps the robot's heading instead of the particle's weights every hypothesis as if it were the odometry.

\item Use the beams that are left after you drop \texttt{inf}; with \(\sigma\)\_z = 0.15 m and \textasciitilde{}110 usable beams a particle inside a wall and one on the correct pose differ by dozens of orders of magnitude, which is why the weights are summed in log form.

\item The scatter you add at the motion step is not decoration: with zero motion noise the cloud can only shrink, and the filter ends up certain about the odometry's own drift.

\item \(\sigma\) from the cloud is the weighted standard deviation around the circular mean of the heading --- the arithmetic mean of +170\textdegree{} and \textminus{}170\textdegree{} points the wrong way.

\end{itemize}

\vspace{4pt}\noindent\textbf{Checkpoint}

\begin{tabular}{|>{\raggedright\hyphenpenalty 500\arraybackslash}p{4.87cm}|>{\raggedright\hyphenpenalty 500\arraybackslash}p{2.80cm}|>{\raggedright\hyphenpenalty 500\arraybackslash}p{4.48cm}|>{\raggedright\hyphenpenalty 500\arraybackslash}p{2.34cm}|}
\hline
\textbf{checkpoint} & \textbf{demo shown} & \textbf{assistant initials} & \textbf{date} \\
\hline
\rule{0pt}{2.4em} & \rule{0pt}{2.4em} & \rule{0pt}{2.4em} & \rule{0pt}{2.4em} \\
\hline
\end{tabular}

\vspace{2pt}\noindent AI tools used (language, debugging, API questions only, never the reference solution)\hrulefill

\newpage

\section{L2 --- Parked somewhere else: localisation from a 2 m prior}

\noindent\hfill\pts{35}\hspace{1em} hall: \texttt{production}\hspace{1em}\texttt{task: mcl\_wide}

\vspace{2pt}\noindent\textbf{Task}

Same hall, same drive, and the robot was not where the odometry claims: the prior around the start pose has a 1\(\sigma\) of two metres instead of half a metre, so the true pose is in the cloud but surrounded by a thousand wrong neighbours. What has to survive is the weighting --- the first scans have to throw away the neighbours. \texttt{max\_error} is the number that grades the beginning of the run rather than its average: a filter that takes eight seconds to find the robot has a good RMSE and a useless first eight seconds.

\vspace{2pt}\noindent\textbf{What is graded}

\begin{tabular}{|>{\raggedright\hyphenpenalty 500\arraybackslash}p{9.20cm}|>{\raggedright\hyphenpenalty 500\arraybackslash}p{6.00cm}|}
\hline
\textbf{quantity} & \textbf{target (from the task file)} \\
\hline
accuracy (RMSE over the graded window) & \(\le\) 60 mm \\
improvement over raw odometry & \(\ge\) 3\(\times\) \\
worst single error & \(\le\) 350 mm \\
rate of kf/pose & \(\ge\) 5 Hz \\
wall contacts & \(\le\) 0 \\
NEES (\(\sigma\) honest) & [0.05--8] \\
\hline
\end{tabular}

\vspace{2pt}\noindent\textbf{Protocol} --- one table per item. Four blank rows each; add rows if you need them, and write the command that produced a number under the table.

\vspace{4pt}\noindent\textbf{C1.2} \emph{The initial cloud is drawn from the prior of the task, not placed on the odometry pose with zero spread}

\begin{tabular}{|>{\raggedright\hyphenpenalty 500\arraybackslash}p{2.87cm}|>{\raggedright\hyphenpenalty 500\arraybackslash}p{2.28cm}|>{\raggedright\hyphenpenalty 500\arraybackslash}p{2.87cm}|>{\raggedright\hyphenpenalty 500\arraybackslash}p{2.28cm}|>{\raggedright\hyphenpenalty 500\arraybackslash}p{3.70cm}|}
\hline
\textbf{question / item} & \textbf{method} & \textbf{measured} & \textbf{target} & \textbf{explanation of the deviation} \\
\hline
\rule{0pt}{2.4em} & \rule{0pt}{2.4em} & \rule{0pt}{2.4em} & \rule{0pt}{2.4em} & \rule{0pt}{2.4em} \\
\rule{0pt}{2.4em} & \rule{0pt}{2.4em} & \rule{0pt}{2.4em} & \rule{0pt}{2.4em} & \rule{0pt}{2.4em} \\
\rule{0pt}{2.4em} & \rule{0pt}{2.4em} & \rule{0pt}{2.4em} & \rule{0pt}{2.4em} & \rule{0pt}{2.4em} \\
\rule{0pt}{2.4em} & \rule{0pt}{2.4em} & \rule{0pt}{2.4em} & \rule{0pt}{2.4em} & \rule{0pt}{2.4em} \\
\hline
\end{tabular}

\vspace{4pt}\noindent\textbf{C2.2} \emph{The first scans reduce N\_eff and then the resampling concentrates the cloud (watch it in RViz/the window)}

\begin{tabular}{|>{\raggedright\hyphenpenalty 500\arraybackslash}p{2.87cm}|>{\raggedright\hyphenpenalty 500\arraybackslash}p{2.28cm}|>{\raggedright\hyphenpenalty 500\arraybackslash}p{2.87cm}|>{\raggedright\hyphenpenalty 500\arraybackslash}p{2.28cm}|>{\raggedright\hyphenpenalty 500\arraybackslash}p{3.70cm}|}
\hline
\textbf{question / item} & \textbf{method} & \textbf{measured} & \textbf{target} & \textbf{explanation of the deviation} \\
\hline
\rule{0pt}{2.4em} & \rule{0pt}{2.4em} & \rule{0pt}{2.4em} & \rule{0pt}{2.4em} & \rule{0pt}{2.4em} \\
\rule{0pt}{2.4em} & \rule{0pt}{2.4em} & \rule{0pt}{2.4em} & \rule{0pt}{2.4em} & \rule{0pt}{2.4em} \\
\rule{0pt}{2.4em} & \rule{0pt}{2.4em} & \rule{0pt}{2.4em} & \rule{0pt}{2.4em} & \rule{0pt}{2.4em} \\
\rule{0pt}{2.4em} & \rule{0pt}{2.4em} & \rule{0pt}{2.4em} & \rule{0pt}{2.4em} & \rule{0pt}{2.4em} \\
\hline
\end{tabular}

\vspace{4pt}\noindent\textbf{C3.2} \emph{No use of /<robot>/truth, and no GPS: this task is decided by the LIDAR and the map}

\begin{tabular}{|>{\raggedright\hyphenpenalty 500\arraybackslash}p{2.87cm}|>{\raggedright\hyphenpenalty 500\arraybackslash}p{2.28cm}|>{\raggedright\hyphenpenalty 500\arraybackslash}p{2.87cm}|>{\raggedright\hyphenpenalty 500\arraybackslash}p{2.28cm}|>{\raggedright\hyphenpenalty 500\arraybackslash}p{3.70cm}|}
\hline
\textbf{question / item} & \textbf{method} & \textbf{measured} & \textbf{target} & \textbf{explanation of the deviation} \\
\hline
\rule{0pt}{2.4em} & \rule{0pt}{2.4em} & \rule{0pt}{2.4em} & \rule{0pt}{2.4em} & \rule{0pt}{2.4em} \\
\rule{0pt}{2.4em} & \rule{0pt}{2.4em} & \rule{0pt}{2.4em} & \rule{0pt}{2.4em} & \rule{0pt}{2.4em} \\
\rule{0pt}{2.4em} & \rule{0pt}{2.4em} & \rule{0pt}{2.4em} & \rule{0pt}{2.4em} & \rule{0pt}{2.4em} \\
\rule{0pt}{2.4em} & \rule{0pt}{2.4em} & \rule{0pt}{2.4em} & \rule{0pt}{2.4em} & \rule{0pt}{2.4em} \\
\hline
\end{tabular}

\vspace{4pt}\noindent\textbf{C4.2} \emph{Report of how many scans the first 25 cm took --- from the counters on kf/info, not from a stopwatch}

\begin{tabular}{|>{\raggedright\hyphenpenalty 500\arraybackslash}p{2.87cm}|>{\raggedright\hyphenpenalty 500\arraybackslash}p{2.28cm}|>{\raggedright\hyphenpenalty 500\arraybackslash}p{2.87cm}|>{\raggedright\hyphenpenalty 500\arraybackslash}p{2.28cm}|>{\raggedright\hyphenpenalty 500\arraybackslash}p{3.70cm}|}
\hline
\textbf{question / item} & \textbf{method} & \textbf{measured} & \textbf{target} & \textbf{explanation of the deviation} \\
\hline
\rule{0pt}{2.4em} & \rule{0pt}{2.4em} & \rule{0pt}{2.4em} & \rule{0pt}{2.4em} & \rule{0pt}{2.4em} \\
\rule{0pt}{2.4em} & \rule{0pt}{2.4em} & \rule{0pt}{2.4em} & \rule{0pt}{2.4em} & \rule{0pt}{2.4em} \\
\rule{0pt}{2.4em} & \rule{0pt}{2.4em} & \rule{0pt}{2.4em} & \rule{0pt}{2.4em} & \rule{0pt}{2.4em} \\
\rule{0pt}{2.4em} & \rule{0pt}{2.4em} & \rule{0pt}{2.4em} & \rule{0pt}{2.4em} & \rule{0pt}{2.4em} \\
\hline
\end{tabular}

\vspace{4pt}\noindent\textbf{Pointers}

\begin{itemize}\itemsep1pt

\item A prior that is too narrow is not optimism, it is a different filter: if no particle is within \(\sigma\)\_z of the true pose, no amount of weighting will find it.

\item Look at N\_eff while it converges. Falling and then rising again is the shape of a filter that found the robot; sitting at 1 is a filter that found one lucky particle instead.

\item If you are tempted to add random particles to make this pass: measure what that costs in L1 first (see docs/mcl.md) --- injection is a recovery mechanism, not a performance knob.

\end{itemize}

\vspace{4pt}\noindent\textbf{Checkpoint}

\begin{tabular}{|>{\raggedright\hyphenpenalty 500\arraybackslash}p{4.87cm}|>{\raggedright\hyphenpenalty 500\arraybackslash}p{2.80cm}|>{\raggedright\hyphenpenalty 500\arraybackslash}p{4.48cm}|>{\raggedright\hyphenpenalty 500\arraybackslash}p{2.34cm}|}
\hline
\textbf{checkpoint} & \textbf{demo shown} & \textbf{assistant initials} & \textbf{date} \\
\hline
\rule{0pt}{2.4em} & \rule{0pt}{2.4em} & \rule{0pt}{2.4em} & \rule{0pt}{2.4em} \\
\hline
\end{tabular}

\vspace{2pt}\noindent AI tools used (language, debugging, API questions only, never the reference solution)\hrulefill

\newpage

\section{L3 --- Same accuracy, a quarter of the particles: where the budget belongs}

\noindent\hfill\pts{35}\hspace{1em} hall: \texttt{production}\hspace{1em}\texttt{task: mcl\_budget}

\vspace{2pt}\noindent\textbf{Task}

The task file now says 250 particles and \textbf{every} beam instead of one in three --- a quarter of the particles for nearly the same accuracy, and 3\(\times\) the beams per particle. Both settings cost the same per scan (particles \(\times\) beams is the product that sets the load), so this is the same computation spent differently, and the accuracy you get for it is not the same: 360 beams constrain a pose far more sharply than 1200 particles cover a hall. The thresholds below are L1's, loosened by as much as a quarter of a cloud measurably costs --- that difference is itself a result worth a line in your protocol. Run \texttt{tools/mcl\_report.py\ --timing} on your own filter and hand in the table of what each knob costs.

\vspace{2pt}\noindent\textbf{What is graded}

\begin{tabular}{|>{\raggedright\hyphenpenalty 500\arraybackslash}p{9.20cm}|>{\raggedright\hyphenpenalty 500\arraybackslash}p{6.00cm}|}
\hline
\textbf{quantity} & \textbf{target (from the task file)} \\
\hline
accuracy (RMSE over the graded window) & \(\le\) 70 mm \\
improvement over raw odometry & \(\ge\) 3\(\times\) \\
worst single error & \(\le\) 300 mm \\
rate of kf/pose & \(\ge\) 5 Hz \\
wall contacts & \(\le\) 0 \\
NEES (\(\sigma\) honest) & [0.05--12] \\
\hline
\end{tabular}

\vspace{2pt}\noindent\textbf{Protocol} --- one table per item. Four blank rows each; add rows if you need them, and write the command that produced a number under the table.

\vspace{4pt}\noindent\textbf{C1.3} \emph{One update: all particles and all beams in vectorised NumPy, no Python loop over particles}

\begin{tabular}{|>{\raggedright\hyphenpenalty 500\arraybackslash}p{2.87cm}|>{\raggedright\hyphenpenalty 500\arraybackslash}p{2.28cm}|>{\raggedright\hyphenpenalty 500\arraybackslash}p{2.87cm}|>{\raggedright\hyphenpenalty 500\arraybackslash}p{2.28cm}|>{\raggedright\hyphenpenalty 500\arraybackslash}p{3.70cm}|}
\hline
\textbf{question / item} & \textbf{method} & \textbf{measured} & \textbf{target} & \textbf{explanation of the deviation} \\
\hline
\rule{0pt}{2.4em} & \rule{0pt}{2.4em} & \rule{0pt}{2.4em} & \rule{0pt}{2.4em} & \rule{0pt}{2.4em} \\
\rule{0pt}{2.4em} & \rule{0pt}{2.4em} & \rule{0pt}{2.4em} & \rule{0pt}{2.4em} & \rule{0pt}{2.4em} \\
\rule{0pt}{2.4em} & \rule{0pt}{2.4em} & \rule{0pt}{2.4em} & \rule{0pt}{2.4em} & \rule{0pt}{2.4em} \\
\rule{0pt}{2.4em} & \rule{0pt}{2.4em} & \rule{0pt}{2.4em} & \rule{0pt}{2.4em} & \rule{0pt}{2.4em} \\
\hline
\end{tabular}

\vspace{4pt}\noindent\textbf{C2.3} \emph{The product particles \(\times\) beams as the cost of one update, measured rather than estimated}

\begin{tabular}{|>{\raggedright\hyphenpenalty 500\arraybackslash}p{2.87cm}|>{\raggedright\hyphenpenalty 500\arraybackslash}p{2.28cm}|>{\raggedright\hyphenpenalty 500\arraybackslash}p{2.87cm}|>{\raggedright\hyphenpenalty 500\arraybackslash}p{2.28cm}|>{\raggedright\hyphenpenalty 500\arraybackslash}p{3.70cm}|}
\hline
\textbf{question / item} & \textbf{method} & \textbf{measured} & \textbf{target} & \textbf{explanation of the deviation} \\
\hline
\rule{0pt}{2.4em} & \rule{0pt}{2.4em} & \rule{0pt}{2.4em} & \rule{0pt}{2.4em} & \rule{0pt}{2.4em} \\
\rule{0pt}{2.4em} & \rule{0pt}{2.4em} & \rule{0pt}{2.4em} & \rule{0pt}{2.4em} & \rule{0pt}{2.4em} \\
\rule{0pt}{2.4em} & \rule{0pt}{2.4em} & \rule{0pt}{2.4em} & \rule{0pt}{2.4em} & \rule{0pt}{2.4em} \\
\rule{0pt}{2.4em} & \rule{0pt}{2.4em} & \rule{0pt}{2.4em} & \rule{0pt}{2.4em} & \rule{0pt}{2.4em} \\
\hline
\end{tabular}

\vspace{4pt}\noindent\textbf{C3.3} \emph{What broke first when the particles were taken away: N\_eff, the RMSE, or the maximum error}

\begin{tabular}{|>{\raggedright\hyphenpenalty 500\arraybackslash}p{4.29cm}|>{\raggedright\hyphenpenalty 500\arraybackslash}p{2.95cm}|>{\raggedright\hyphenpenalty 500\arraybackslash}p{4.29cm}|>{\raggedright\hyphenpenalty 500\arraybackslash}p{2.95cm}|}
\hline
\textbf{quantity} & \textbf{value} & \textbf{how it was measured} & \textbf{what could make it wrong} \\
\hline
\rule{0pt}{2.4em} & \rule{0pt}{2.4em} & \rule{0pt}{2.4em} & \rule{0pt}{2.4em} \\
\rule{0pt}{2.4em} & \rule{0pt}{2.4em} & \rule{0pt}{2.4em} & \rule{0pt}{2.4em} \\
\rule{0pt}{2.4em} & \rule{0pt}{2.4em} & \rule{0pt}{2.4em} & \rule{0pt}{2.4em} \\
\rule{0pt}{2.4em} & \rule{0pt}{2.4em} & \rule{0pt}{2.4em} & \rule{0pt}{2.4em} \\
\hline
\end{tabular}

\vspace{4pt}\noindent\textbf{C4.3} \emph{A short answer to why 360 beams are not 360 independent measurements of a hall}

\begin{tabular}{|>{\raggedright\hyphenpenalty 500\arraybackslash}p{2.87cm}|>{\raggedright\hyphenpenalty 500\arraybackslash}p{2.28cm}|>{\raggedright\hyphenpenalty 500\arraybackslash}p{2.87cm}|>{\raggedright\hyphenpenalty 500\arraybackslash}p{2.28cm}|>{\raggedright\hyphenpenalty 500\arraybackslash}p{3.70cm}|}
\hline
\textbf{question / item} & \textbf{method} & \textbf{measured} & \textbf{target} & \textbf{explanation of the deviation} \\
\hline
\rule{0pt}{2.4em} & \rule{0pt}{2.4em} & \rule{0pt}{2.4em} & \rule{0pt}{2.4em} & \rule{0pt}{2.4em} \\
\rule{0pt}{2.4em} & \rule{0pt}{2.4em} & \rule{0pt}{2.4em} & \rule{0pt}{2.4em} & \rule{0pt}{2.4em} \\
\rule{0pt}{2.4em} & \rule{0pt}{2.4em} & \rule{0pt}{2.4em} & \rule{0pt}{2.4em} & \rule{0pt}{2.4em} \\
\rule{0pt}{2.4em} & \rule{0pt}{2.4em} & \rule{0pt}{2.4em} & \rule{0pt}{2.4em} & \rule{0pt}{2.4em} \\
\hline
\end{tabular}

\vspace{4pt}\noindent\textbf{Pointers}

\begin{itemize}\itemsep1pt

\item The work of one update is (particles \(\times\) beams) gathers in the distance field. Halving one factor and tripling the other keeps the work the same and changes what the filter is good at.

\item Fewer particles means a coarser sampling of the pose space: \(\sigma\)\_z controls how forgiving the likelihood is about that coarseness, and it is the knob that lets 250 particles find a 1 cm peak.

\item Report N\_eff for both settings. If it did not go down with a quarter of the cloud, the resampling is doing something you did not intend.

\end{itemize}

\vspace{4pt}\noindent\textbf{Checkpoint}

\begin{tabular}{|>{\raggedright\hyphenpenalty 500\arraybackslash}p{4.87cm}|>{\raggedright\hyphenpenalty 500\arraybackslash}p{2.80cm}|>{\raggedright\hyphenpenalty 500\arraybackslash}p{4.48cm}|>{\raggedright\hyphenpenalty 500\arraybackslash}p{2.34cm}|}
\hline
\textbf{checkpoint} & \textbf{demo shown} & \textbf{assistant initials} & \textbf{date} \\
\hline
\rule{0pt}{2.4em} & \rule{0pt}{2.4em} & \rule{0pt}{2.4em} & \rule{0pt}{2.4em} \\
\hline
\end{tabular}

\vspace{2pt}\noindent AI tools used (language, debugging, API questions only, never the reference solution)\hrulefill

\newpage

\section{L4 --- The window is dirty: make the model match the sensor}

\noindent\hfill\pts{30}\hspace{1em} hall: \texttt{production}\hspace{1em}\texttt{task: mcl\_dirty}

\vspace{2pt}\noindent\textbf{Task}

The same hall, the same drive and the same filter as L1, with one change you will not see in the topic: the LIDAR's \(\sigma\) is 0.25 m instead of 0.015 m --- a reflective, partially covered window, which on a real robot is a Tuesday. The parameters that won L1 were tuned against a clean sensor and they will not survive this one. Not because the filter is broken, but because \(\sigma\)\_z is not a knob on the filter, it is a claim about the sensor: weight a 250 mm measurement as if it were a 150 mm measurement and the cloud collapses onto whatever the first dirty beams liked, and the reported \(\sigma\) becomes a lie even while the position happens to stay reasonable. Find parameters that work with this sensor and write down which one you changed and why that one.

\vspace{2pt}\noindent\textbf{What is graded}

\begin{tabular}{|>{\raggedright\hyphenpenalty 500\arraybackslash}p{9.20cm}|>{\raggedright\hyphenpenalty 500\arraybackslash}p{6.00cm}|}
\hline
\textbf{quantity} & \textbf{target (from the task file)} \\
\hline
accuracy (RMSE over the graded window) & \(\le\) 70 mm \\
improvement over raw odometry & \(\ge\) 2.5\(\times\) \\
worst single error & \(\le\) 200 mm \\
rate of kf/pose & \(\ge\) 5 Hz \\
wall contacts & \(\le\) 0 \\
NEES (\(\sigma\) honest) & [0.3--5] \\
\hline
\end{tabular}

\vspace{2pt}\noindent\textbf{Protocol} --- one table per item. Four blank rows each; add rows if you need them, and write the command that produced a number under the table.

\vspace{4pt}\noindent\textbf{C1.4} \emph{RMSE and improvement over the raw odometry, with the parameters that achieved them}

\begin{tabular}{|>{\raggedright\hyphenpenalty 500\arraybackslash}p{2.75cm}|>{\raggedright\hyphenpenalty 500\arraybackslash}p{1.58cm}|>{\raggedright\hyphenpenalty 500\arraybackslash}p{1.58cm}|>{\raggedright\hyphenpenalty 500\arraybackslash}p{3.55cm}|>{\raggedright\hyphenpenalty 500\arraybackslash}p{1.58cm}|>{\raggedright\hyphenpenalty 500\arraybackslash}p{2.47cm}|}
\hline
\textbf{run / settings} & \textbf{RMSE [mm]} & \textbf{raw odom [mm]} & \textbf{improvement [-]} & \textbf{NEES [-]} & \textbf{verdict} \\
\hline
\rule{0pt}{2.4em} & \rule{0pt}{2.4em} & \rule{0pt}{2.4em} & \rule{0pt}{2.4em} & \rule{0pt}{2.4em} & \rule{0pt}{2.4em} \\
\rule{0pt}{2.4em} & \rule{0pt}{2.4em} & \rule{0pt}{2.4em} & \rule{0pt}{2.4em} & \rule{0pt}{2.4em} & \rule{0pt}{2.4em} \\
\rule{0pt}{2.4em} & \rule{0pt}{2.4em} & \rule{0pt}{2.4em} & \rule{0pt}{2.4em} & \rule{0pt}{2.4em} & \rule{0pt}{2.4em} \\
\rule{0pt}{2.4em} & \rule{0pt}{2.4em} & \rule{0pt}{2.4em} & \rule{0pt}{2.4em} & \rule{0pt}{2.4em} & \rule{0pt}{2.4em} \\
\hline
\end{tabular}

\vspace{4pt}\noindent\textbf{C2.4} \emph{NEES: does your \(\sigma\) describe your error now, and what happened to it when you changed \(\sigma\)\_z}

\begin{tabular}{|>{\raggedright\hyphenpenalty 500\arraybackslash}p{4.45cm}|>{\raggedright\hyphenpenalty 500\arraybackslash}p{3.02cm}|>{\raggedright\hyphenpenalty 500\arraybackslash}p{2.40cm}|>{\raggedright\hyphenpenalty 500\arraybackslash}p{1.73cm}|>{\raggedright\hyphenpenalty 500\arraybackslash}p{2.40cm}|}
\hline
\textbf{configuration} & \textbf{reported \(\sigma\) [mm]} & \textbf{actual error [mm]} & \textbf{NEES [-]} & \textbf{should be [-]} \\
\hline
\rule{0pt}{2.4em} & \rule{0pt}{2.4em} & \rule{0pt}{2.4em} & \rule{0pt}{2.4em} & \rule{0pt}{2.4em} \\
\rule{0pt}{2.4em} & \rule{0pt}{2.4em} & \rule{0pt}{2.4em} & \rule{0pt}{2.4em} & \rule{0pt}{2.4em} \\
\rule{0pt}{2.4em} & \rule{0pt}{2.4em} & \rule{0pt}{2.4em} & \rule{0pt}{2.4em} & \rule{0pt}{2.4em} \\
\rule{0pt}{2.4em} & \rule{0pt}{2.4em} & \rule{0pt}{2.4em} & \rule{0pt}{2.4em} & \rule{0pt}{2.4em} \\
\hline
\end{tabular}

\vspace{4pt}\noindent\textbf{C3.4} \emph{A \(\sigma\)\_z sweep --- at least four values, one number each --- as a table in the protocol}

\begin{tabular}{|>{\raggedright\hyphenpenalty 500\arraybackslash}p{2.90cm}|>{\raggedright\hyphenpenalty 500\arraybackslash}p{1.81cm}|>{\raggedright\hyphenpenalty 500\arraybackslash}p{1.52cm}|>{\raggedright\hyphenpenalty 500\arraybackslash}p{3.40cm}|>{\raggedright\hyphenpenalty 500\arraybackslash}p{1.52cm}|>{\raggedright\hyphenpenalty 500\arraybackslash}p{2.37cm}|}
\hline
\textbf{parameter} & \textbf{value tried} & \textbf{RMSE [mm]} & \textbf{improvement [-]} & \textbf{NEES [-]} & \textbf{comment} \\
\hline
\rule{0pt}{2.4em} & \rule{0pt}{2.4em} & \rule{0pt}{2.4em} & \rule{0pt}{2.4em} & \rule{0pt}{2.4em} & \rule{0pt}{2.4em} \\
\rule{0pt}{2.4em} & \rule{0pt}{2.4em} & \rule{0pt}{2.4em} & \rule{0pt}{2.4em} & \rule{0pt}{2.4em} & \rule{0pt}{2.4em} \\
\rule{0pt}{2.4em} & \rule{0pt}{2.4em} & \rule{0pt}{2.4em} & \rule{0pt}{2.4em} & \rule{0pt}{2.4em} & \rule{0pt}{2.4em} \\
\rule{0pt}{2.4em} & \rule{0pt}{2.4em} & \rule{0pt}{2.4em} & \rule{0pt}{2.4em} & \rule{0pt}{2.4em} & \rule{0pt}{2.4em} \\
\hline
\end{tabular}

\vspace{4pt}\noindent\textbf{C4.4} \emph{Why \(\sigma\)\_z \(\approx\) 0.5 m works better here than \(\sigma\)\_z = 0.25 m, which is the \(\sigma\) the sensor actually has}

\begin{tabular}{|>{\raggedright\hyphenpenalty 500\arraybackslash}p{2.87cm}|>{\raggedright\hyphenpenalty 500\arraybackslash}p{2.28cm}|>{\raggedright\hyphenpenalty 500\arraybackslash}p{2.87cm}|>{\raggedright\hyphenpenalty 500\arraybackslash}p{2.28cm}|>{\raggedright\hyphenpenalty 500\arraybackslash}p{3.70cm}|}
\hline
\textbf{question / item} & \textbf{method} & \textbf{measured} & \textbf{target} & \textbf{explanation of the deviation} \\
\hline
\rule{0pt}{2.4em} & \rule{0pt}{2.4em} & \rule{0pt}{2.4em} & \rule{0pt}{2.4em} & \rule{0pt}{2.4em} \\
\rule{0pt}{2.4em} & \rule{0pt}{2.4em} & \rule{0pt}{2.4em} & \rule{0pt}{2.4em} & \rule{0pt}{2.4em} \\
\rule{0pt}{2.4em} & \rule{0pt}{2.4em} & \rule{0pt}{2.4em} & \rule{0pt}{2.4em} & \rule{0pt}{2.4em} \\
\rule{0pt}{2.4em} & \rule{0pt}{2.4em} & \rule{0pt}{2.4em} & \rule{0pt}{2.4em} & \rule{0pt}{2.4em} \\
\hline
\end{tabular}

\vspace{4pt}\noindent\textbf{C5.4} \emph{What \texttt{z\_rand} costs you when the sensor is this noisy}

\begin{tabular}{|>{\raggedright\hyphenpenalty 500\arraybackslash}p{4.29cm}|>{\raggedright\hyphenpenalty 500\arraybackslash}p{2.95cm}|>{\raggedright\hyphenpenalty 500\arraybackslash}p{4.29cm}|>{\raggedright\hyphenpenalty 500\arraybackslash}p{2.95cm}|}
\hline
\textbf{quantity} & \textbf{value} & \textbf{how it was measured} & \textbf{what could make it wrong} \\
\hline
\rule{0pt}{2.4em} & \rule{0pt}{2.4em} & \rule{0pt}{2.4em} & \rule{0pt}{2.4em} \\
\rule{0pt}{2.4em} & \rule{0pt}{2.4em} & \rule{0pt}{2.4em} & \rule{0pt}{2.4em} \\
\rule{0pt}{2.4em} & \rule{0pt}{2.4em} & \rule{0pt}{2.4em} & \rule{0pt}{2.4em} \\
\rule{0pt}{2.4em} & \rule{0pt}{2.4em} & \rule{0pt}{2.4em} & \rule{0pt}{2.4em} \\
\hline
\end{tabular}

\vspace{4pt}\noindent\textbf{Pointers}

\begin{itemize}\itemsep1pt

\item The measurement noise is in the task's \texttt{sim} block, and the grader checks that the run really used it: \texttt{lidar.sigma} is 0.25 m. Read it off the config rather than assuming 0.015 m --- a filter that asks the bus what its sensor can do is a filter that survives a sensor change.

\item \(\sigma\)\_z enters the weight of every beam, so it is the parameter this task is about. Everything else can stay where L1 left it.

\item Measured on this recording, for comparison once you have your own numbers: \(\sigma\)\_z 0.15 \(\rightarrow\) 96 mm and NEES \textasciitilde{}15; \(\sigma\)\_z 0.30 \(\rightarrow\) 62 mm; \(\sigma\)\_z 0.50 \(\rightarrow\) 49 mm and NEES 2.2; \(\sigma\)\_z 0.80 \(\rightarrow\) 62 mm. There is a minimum in there, and it is near 2\(\cdot\)\(\sigma\)\_sensor rather than at \(\sigma\)\_sensor.

\item \texttt{tools/mcl\_report.py} replays a recording in seconds, so a sweep of four \(\sigma\)\_z values is a loop, not four lab sessions.

\end{itemize}

\vspace{4pt}\noindent\textbf{Checkpoint}

\begin{tabular}{|>{\raggedright\hyphenpenalty 500\arraybackslash}p{4.87cm}|>{\raggedright\hyphenpenalty 500\arraybackslash}p{2.80cm}|>{\raggedright\hyphenpenalty 500\arraybackslash}p{4.48cm}|>{\raggedright\hyphenpenalty 500\arraybackslash}p{2.34cm}|}
\hline
\textbf{checkpoint} & \textbf{demo shown} & \textbf{assistant initials} & \textbf{date} \\
\hline
\rule{0pt}{2.4em} & \rule{0pt}{2.4em} & \rule{0pt}{2.4em} & \rule{0pt}{2.4em} \\
\hline
\end{tabular}

\vspace{2pt}\noindent AI tools used (language, debugging, API questions only, never the reference solution)\hrulefill

\newpage

\end{document}
